Dans l'espace muni d'un repère orthonormal \Oijk, on considère les points 

$A(-1~;~1~;~1)$, $B(1~;~-6~;~-1)$ et $C(-5~;~2~;~3)$.

\medskip

\begin{enumerate}
\item[1.] 
	\begin{enumerate}[label=\alph*.]
		\item Démontrer que les points A, B et C ne sont pas alignés.
		\item Démontrer que le vecteur $\vec{n}\begin{pmatrix} 6 \\ -2 \\ 13 \end{pmatrix}$ est normal au plan (ABC).
		\item En déduire que le plan (ABC) admet pour équation cartésienne :

\[6x - 2y + 13z - 5 = 0.\]
	\end{enumerate}

\item On considère le plan $(P)$ d'équation cartésienne $8x - 2y - 4z + 18 = 0$.

On admet que deux plans sont perpendiculaires si leurs vecteurs normaux respectifs sont orthogonaux.

Démontrer que les plans (ABC) et $(P)$ sont perpendiculaires.
\item Soit $(d)$ une droite dont une représentation paramétrique est :

\[\begin{cases}
x = \phantom{1}-4 + \phantom{1}t \\
y = -11 + 6t \\
z = \phantom{-1}2 - \phantom{1}t
\end{cases}, \quad t \in \R.\]

Démontrer que la droite $(d)$ intersecte le plan (ABC) en un unique point qu'on nommera E, et qui a pour coordonnées E$(-3~;~-5~;~1)$.

\item[4.] Soit H$(-2~;~1~;~0)$.
	\begin{enumerate}[label=\alph*.]
		\item Montrer que H est le projeté orthogonal de A sur $(d)$.
		\item En déduire la distance du point A à la droite $(d)$.
		\item Montrer que le triangle AHE est rectangle et que son aire est égale à $\sqrt{19}$ unités d'aire.
	\end{enumerate}
\end{enumerate}

\bigskip
